\documentclass[11pt,a4paper]{article}

\usepackage[spanish]{babel}
\usepackage[utf8]{inputenc}
\usepackage[T1]{fontenc}
\usepackage{lmodern}
\usepackage{geometry}
\usepackage{xcolor}
\usepackage{listings}
\usepackage{enumitem}
\usepackage{hyperref}

\geometry{margin=2.5cm}

\lstdefinestyle{cstyle}{
  language=C++,
  basicstyle=\ttfamily\small,
  keywordstyle=\bfseries\color{blue!60!black},
  commentstyle=\itshape\color{green!40!black},
  stringstyle=\color{red!50!black},
  numbers=left,
  numberstyle=\tiny,
  stepnumber=1,
  numbersep=8pt,
  showstringspaces=false,
  tabsize=2,
  breaklines=true,
  frame=single,
  columns=fullflexible
}

\lstset{style=cstyle}

\title{Práctica de laboratorio: Sobrecarga de operadores en C++}
\author{Programación II}
\date{}

\begin{document}

\maketitle

\section{Presentación}

En esta práctica implementará una clase \lstinline|Fraction| para representar fracciones racionales. El objetivo es estudiar la sobrecarga de operadores en C++ y distinguir cuándo conviene implementar un operador como \textbf{función miembro} y cuándo conviene implementarlo como \textbf{función no miembro}.

La práctica también reforzará el uso correcto de constructores, encapsulación, métodos \lstinline|const|, manejo básico de errores de entrada y mantenimiento de invariantes dentro de una clase.

\section{Objetivos}

Al finalizar la práctica, el estudiante será capaz de:

\begin{enumerate}
  \item Implementar una clase en C++ con atributos privados y constructores correctos.
  \item Mantener invariantes internos de una clase.
  \item Sobrecargar operadores aritméticos y de comparación.
  \item Distinguir entre operadores implementados como funciones miembro y operadores implementados como funciones no miembro.
  \item Implementar correctamente los operadores de entrada y salida \lstinline|operator>>| y \lstinline|operator<<|.
  \item Usar pruebas simples para verificar el comportamiento de una clase.
\end{enumerate}

\section{Contexto del problema}

Se desea implementar una clase \lstinline|Fraction| que permita trabajar con fracciones racionales de la forma \lstinline|numerador/denominador|.

Por ejemplo:

\begin{lstlisting}
Fraction a(1, 2);
Fraction b(1, 3);

Fraction c = a + b;

std::cout << c << '\n';
\end{lstlisting}

La salida esperada sería:

\begin{lstlisting}[language={}]
5/6
\end{lstlisting}

La clase debe garantizar que toda fracción se mantenga siempre en un estado válido. En particular, el denominador nunca debe ser cero.

\section{Archivos del proyecto}

El código base del laboratorio estará organizado de la siguiente forma:

\begin{lstlisting}[language={}]
.
|-- Makefile
|-- README.md
|-- include
|   `-- Fraction.h
|-- src
|   `-- Fraction.cc
`-- tests
    `-- main.cc
\end{lstlisting}

El archivo \lstinline|Fraction.h| contendrá la declaración de la clase y de las funciones asociadas. El archivo \lstinline|Fraction.cc| contendrá la implementación. El archivo \lstinline|tests/main.cc| contendrá pruebas simples para verificar el funcionamiento del código.

\section{Requisitos de implementación}

\subsection{Clase \texttt{Fraction}}

Debe implementar una clase \lstinline|Fraction| con dos atributos privados:

\begin{lstlisting}
int numerator;
int denominator;
\end{lstlisting}

La clase debe ofrecer, como mínimo, los siguientes constructores y métodos:

\begin{lstlisting}
class Fraction {
private:
  int numerator;
  int denominator;

public:
  Fraction(int numerator, int denominator);
  explicit Fraction(int numerator);

  int num() const;
  int den() const;
};
\end{lstlisting}

El constructor \lstinline|Fraction(int numerator, int denominator)| debe rechazar denominadores iguales a cero.

El constructor \lstinline|explicit Fraction(int numerator)| debe construir una fracción entera, equivalente a \lstinline|numerator/1|.

Los métodos \lstinline|num()| y \lstinline|den()| deben ser métodos \lstinline|const|.

\section{Invariante de la clase}

La clase \lstinline|Fraction| debe mantener el siguiente invariante:

\begin{center}
\lstinline|denominator != 0|
\end{center}

Opcionalmente, puede mantener también las fracciones normalizadas. Por ejemplo:

\begin{lstlisting}[language={}]
2/4  ->  1/2
-2/4 -> -1/2
2/-4 -> -1/2
\end{lstlisting}

Si decide implementar normalización, asegúrese de que el signo negativo quede en el numerador y que el denominador sea siempre positivo.

\section{Sobrecarga de operadores como funciones miembro}

Algunos operadores modifican directamente el objeto ubicado a la izquierda del operador. En estos casos, es natural implementarlos como \textbf{funciones miembro}.

Debe implementar los siguientes operadores como funciones miembro:

\begin{lstlisting}
Fraction& operator+=(const Fraction& other);
Fraction& operator-=(const Fraction& other);
\end{lstlisting}

Por ejemplo:

\begin{lstlisting}
Fraction a(1, 2);
Fraction b(1, 3);

a += b;
\end{lstlisting}

Después de ejecutar el código anterior, \lstinline|a| debe representar \lstinline|5/6|.

La expresión \lstinline|a += b| se interpreta conceptualmente como \lstinline|a.operator+=(b)|. Por esta razón, \lstinline|operator+=| se implementa naturalmente como función miembro: el objeto izquierdo es el que se modifica.

Ambos operadores deben retornar \lstinline|*this| para permitir expresiones encadenadas.

\section{Sobrecarga de operadores como funciones no miembro}

Otros operadores no modifican ninguno de sus operandos. En estos casos, normalmente conviene implementarlos como \textbf{funciones no miembro}.

Debe implementar los siguientes operadores como funciones no miembro:

\begin{lstlisting}
Fraction operator+(Fraction a, const Fraction& b);
Fraction operator-(Fraction a, const Fraction& b);
Fraction operator*(const Fraction& a, const Fraction& b);
Fraction operator/(const Fraction& a, const Fraction& b);
\end{lstlisting}

La implementación de \lstinline|operator+| puede reutilizar \lstinline|operator+=|:

\begin{lstlisting}
Fraction operator+(Fraction a, const Fraction& b){
  a += b;
  return a;
}
\end{lstlisting}

De manera similar, \lstinline|operator-| puede reutilizar \lstinline|operator-=|.

\section{Operadores de comparación}

Debe implementar, como mínimo, los siguientes operadores de comparación como funciones no miembro:

\begin{lstlisting}
bool operator==(const Fraction& a, const Fraction& b);
bool operator!=(const Fraction& a, const Fraction& b);
bool operator<(const Fraction& a, const Fraction& b);
bool operator<=(const Fraction& a, const Fraction& b);
bool operator>(const Fraction& a, const Fraction& b);
bool operator>=(const Fraction& a, const Fraction& b);
\end{lstlisting}

Sugerencia: no es necesario implementar todos desde cero. Puede definir algunos operadores a partir de \lstinline|operator<| y \lstinline|operator==|.

Para comparar fracciones sin usar números reales, puede comparar mediante productos cruzados: \lstinline|a/b < c/d| si y solo si \lstinline|a*d < c*b|.

\section{Operador de salida \texorpdfstring{\texttt{operator\textless\textless}}{operator<<}}

Debe implementar el operador de salida como función no miembro:

\begin{lstlisting}
std::ostream& operator<<(std::ostream& os, const Fraction& f);
\end{lstlisting}

Este operador debe permitir escribir:

\begin{lstlisting}
Fraction f(3, 4);

std::cout << f << '\n';
\end{lstlisting}

La salida debe ser \lstinline|3/4|.

El operador debe retornar el flujo recibido para permitir encadenar salidas:

\begin{lstlisting}
std::cout << "f = " << f << '\n';
\end{lstlisting}

\lstinline|operator<<| se implementa como función no miembro porque el operando izquierdo es un objeto de tipo \lstinline|std::ostream|, por ejemplo \lstinline|std::cout|.

\section{Operador de entrada \texorpdfstring{\texttt{operator\textgreater\textgreater}}{operator>>}}

Debe implementar el operador de entrada como función no miembro:

\begin{lstlisting}
std::istream& operator>>(std::istream& is, Fraction& f);
\end{lstlisting}

Este operador debe leer fracciones con el formato \lstinline|numerador/denominador|.

Ejemplos de entrada válida:

\begin{lstlisting}[language={}]
3/4
-2/5
10/1
\end{lstlisting}

Si la entrada es válida, \lstinline|f| debe actualizarse con la fracción leída. Si la entrada es inválida, el operador debe marcar el flujo en estado de error usando \lstinline|is.setstate(std::ios::failbit)|.

Casos inválidos:

\begin{lstlisting}[language={}]
3/0
3-4
abc
\end{lstlisting}

El operador no debe dejar al objeto \lstinline|f| en un estado inválido. Sugerencia: primero lea los datos en variables locales. Solo construya y asigne una nueva fracción si la entrada es válida.

\section{Preguntas conceptuales}

Responda brevemente en el archivo \lstinline|README.md| del repositorio:

\begin{enumerate}
  \item ¿Por qué \lstinline|operator+=| se implementa naturalmente como función miembro?
  \item ¿Por qué \lstinline|operator+| puede implementarse como función no miembro?
  \item ¿Por qué \lstinline|operator<<| debe implementarse como función no miembro?
  \item ¿Qué significa que un método sea \lstinline|const|?
  \item ¿Qué invariante debe mantener la clase \lstinline|Fraction|?
  \item ¿Qué debe ocurrir si se intenta construir una fracción con denominador cero?
  \item ¿Por qué \lstinline|operator>>| no debe modificar el objeto si la entrada es inválida?
\end{enumerate}

\section{Pruebas mínimas}

Debe incluir pruebas que verifiquen, al menos, construcción, suma, resta, multiplicación, división, comparación, salida, entrada válida y entrada inválida.

\section{Restricciones}

\begin{enumerate}
  \item Los atributos de \lstinline|Fraction| deben ser privados.
  \item No debe usar variables globales.
  \item No debe representar la fracción usando \lstinline|double| o \lstinline|float|.
  \item No debe implementar operadores aritméticos modificando ambos operandos.
  \item Los operadores \lstinline|<<| y \lstinline|>>| deben retornar referencias al flujo recibido.
  \item Los métodos que no modifican el objeto deben declararse como \lstinline|const|.
  \item La clase debe impedir denominadores iguales a cero.
  \item El código debe compilar sin warnings importantes usando las opciones indicadas en el \lstinline|Makefile|.
\end{enumerate}

\section{Comandos esperados}

El repositorio debe permitir compilar con:

\begin{lstlisting}[language=bash]
make
\end{lstlisting}

Ejecutar las pruebas con:

\begin{lstlisting}[language=bash]
make run
\end{lstlisting}

Limpiar archivos generados con:

\begin{lstlisting}[language=bash]
make clean
\end{lstlisting}

\section{Entregable}

Debe entregar un repositorio con la estructura indicada en esta práctica. El archivo \lstinline|README.md| debe incluir las respuestas a las preguntas conceptuales.

\section{Criterios de evaluación}

Se evaluará:

\begin{enumerate}
  \item Correcta implementación de la clase \lstinline|Fraction|.
  \item Uso adecuado de constructores.
  \item Mantenimiento del invariante \lstinline|denominator != 0|.
  \item Uso correcto de funciones miembro y funciones no miembro para la sobrecarga de operadores.
  \item Correcta implementación de \lstinline|operator<<| y \lstinline|operator>>|.
  \item Uso adecuado de \lstinline|const|.
  \item Claridad y orden del código.
  \item Correcta compilación mediante \lstinline|make|.
  \item Pruebas suficientes en \lstinline|tests/main.cc|.
  \item Respuestas claras en el \lstinline|README.md|.
\end{enumerate}

\section{Reflexión final}

La sobrecarga de operadores debe usarse para que el código sea más natural y expresivo, no para ocultar comportamientos inesperados.

Una clase bien diseñada debe proteger sus invariantes y ofrecer operaciones coherentes con el significado del tipo que representa.

\end{document}
